\paragraph{Toeplitz 主子式的离散凹锥、Blaschke 原点模维数证书、Rouch\'e 边界裕度与 window-\(6\) parity 的几何压缩}
在 Toeplitz--Verblunsky 乘积分解、离切片 Blaschke 指标、Jensen 缺陷的视界纯度判据以及 window-\(6\) boundary uplift 的二层/三层分岔都已明确之后，还可继续抽取出下列五条后果。它们分别把 Toeplitz 主子式链在对数坐标下精确识别为离散凹锥，把单标量 \(|B_{\mathcal Z}(0)|\) 升格为残差维数的 sharp 证书，把 Hurwitz 型零自由逼近进一步有限化为一列圆周上的 Rouch\'e 裕度不等式，并把 window-\(6\) 的 \((\ZZ_2)^3\) boundary parity 先压缩为局域几何只保留的对角 \(\ZZ_2\) 方向，再识别为不能在“无额外协议门、无额外寄存器”的口径下跨 \(m=7,8\) 自然延拓的最小锚点现象。

\begin{theorem}[Toeplitz 主子式锥在对数坐标下恰为离散凹锥]\label{thm:xi-time-part65e-toeplitz-local-logconcavity-realizability}
\leanverified{paper\_xi\_time\_part65e\_toeplitz\_local\_logconcavity\_realizability}
设 \(N\ge 0\)，并给定正数列
\[
\Delta_{-1}:=1,
\qquad
\Delta_0>0,\ \Delta_1>0,\dots,\Delta_N>0.
\]
定义
\[
u_k:=\log \Delta_k
\qquad(-1\le k\le N).
\]
则以下两命题等价：
\begin{enumerate}
  \item 存在某个 Carath\'eodory 函数的截断 Toeplitz 主子式链
  \[
  \mathbf T_k=(c_{j-\ell})_{0\le j,\ell\le k}
  \qquad(0\le k\le N)
  \]
  使得
  \[
  \det \mathbf T_k=\Delta_k
  \qquad(0\le k\le N);
  \]
  \item 对所有 \(0\le k\le N-1\)，都有
  \[
  u_{k+1}-2u_k+u_{k-1}\le 0.
  \]
\end{enumerate}
并且在该情形下，对每个 \(0\le k\le N-1\) 都有严格恒等式
\[
\boxed{\;
u_{k+1}-2u_k+u_{k-1}
=
\log(1-|\alpha_k|^2)\le 0,
\;}
\]
从而相应 Verblunsky 系数的模长由主子式序列唯一恢复为
\[
|\alpha_k|^2
=
1-\exp\!\bigl(u_{k+1}-2u_k+u_{k-1}\bigr)
\qquad(0\le k\le N-1).
\]
换言之，Toeplitz--PSD 主子式链在忘却相位后恰等价于离散凹锥条件。
\end{theorem}
\begin{proof}
当 \(N=0\) 时结论平凡：取常值 Carath\'eodory 函数 \(\mathscr C\equiv \Delta_0\) 即可。以下设 \(N\ge 1\)。

若存在所述 Toeplitz 主子式链，则由定理 \ref{thm:xi-toeplitz-principal-minor-discrete-curvature-recovery}，对 \(k=0\) 有
\[
1-|\alpha_0|^2
=
\frac{\Delta_1}{\Delta_0^2}
=
\exp(u_1-2u_0+u_{-1}),
\]
而对 \(1\le k\le N-1\) 有
\[
1-|\alpha_k|^2
=
\frac{\Delta_{k+1}\Delta_{k-1}}{\Delta_k^2}
=
\exp(u_{k+1}-2u_k+u_{k-1}).
\]
两侧取对数即得
\[
u_{k+1}-2u_k+u_{k-1}
=
\log(1-|\alpha_k|^2)\le 0,
\]
并同时得到模长恢复公式。

反过来，若已知
\[
u_{k+1}-2u_k+u_{k-1}\le 0
\qquad(0\le k\le N-1),
\]
则可定义
\[
\rho_k:=\exp(u_{k+1}-2u_k+u_{k-1})\in(0,1],
\qquad
r_k:=\rho_k,
\qquad
|\alpha_k|:=\sqrt{1-\rho_k}\in[0,1)
\qquad(0\le k\le N-1),
\]
并任选相位；例如取全部 \(\alpha_k\in[0,1)\subset\RR\)。再令
\[
\alpha_j:=0
\qquad(j\ge N).
\]
由 Schur--Carath\'eodory 对应，存在某个 Carath\'eodory 函数 \(\mathscr C\) 具有这列 Verblunsky 系数。记其 Toeplitz 主子式为
\[
\widetilde\Delta_k:=\det \widetilde{\mathbf T}_k
\qquad(k\ge 0).
\]
由定理 \ref{thm:xi-toeplitz-principal-minor-discrete-curvature-recovery}，有
\[
\widetilde\Delta_0=\Delta_0,
\qquad
\widetilde\Delta_1
=
\Delta_0^2(1-|\alpha_0|^2)
=
\Delta_0^2\rho_0
=
\Delta_1.
\]
对每个 \(1\le k\le N-1\)，同一定理还给出递推
\[
\widetilde\Delta_{k+1}
=(1-|\alpha_k|^2)\frac{\widetilde\Delta_k^2}{\widetilde\Delta_{k-1}}
=r_k\frac{\widetilde\Delta_k^2}{\widetilde\Delta_{k-1}}.
\]
若已知 \(\widetilde\Delta_{k-1}=\Delta_{k-1}\) 与 \(\widetilde\Delta_k=\Delta_k\)，则
\[
\widetilde\Delta_{k+1}
=r_k\frac{\Delta_k^2}{\Delta_{k-1}}
=\Delta_{k+1}.
\]
归纳即得 \(\widetilde\Delta_k=\Delta_k\) 对所有 \(0\le k\le N\) 成立，故所需实现性得证。最后，模长唯一恢复公式已由上面的曲率恒等式直接给出。证毕。
\end{proof}

\begin{theorem}[Blaschke 原点模对离切片残差维数的 sharp 上界]\label{thm:xi-time-part65e-blaschke-origin-modulus-dimension-bound}
\leanverified{paper\_xi\_time\_part65e\_blaschke\_origin\_modulus\_dimension\_bound}
沿用定义 \ref{def:xi-offslice-ind-weight} 的记号，并设 \(\mathcal Z_{+}\neq\varnothing\) 为有限离切片零点多重集，\(B_{\mathcal Z}\) 为其对应的有限 Blaschke 乘积。定义
\[
\delta_{\min}(\mathcal Z)
:=
\min_{s\in\mathcal Z_{+}}
\Bigl(\Re(s)-\frac12\Bigr)
>0.
\]
则有 sharp 估计
\[
\mathrm{Ind}(\mathcal Z)
\le
\frac{\mathrm W(\mathcal Z)}{2L\,\delta_{\min}(\mathcal Z)}
=
\frac{-\log|B_{\mathcal Z}(0)|}{2L\,\delta_{\min}(\mathcal Z)}.
\]
等价地，
\[
\dim K_{B_{\mathcal Z}}
\le
\frac{-\log|B_{\mathcal Z}(0)|}{2L\,\delta_{\min}(\mathcal Z)}.
\]
并且等号成立当且仅当全部离切片零点具有相同的实部偏差，即
\[
\Re(s)-\frac12\equiv \delta_{\min}(\mathcal Z)
\qquad(s\in\mathcal Z_{+}).
\]
\end{theorem}
\begin{proof}
由定理 \ref{thm:xi-offslice-realpart-sum-law}，
\[
\mathrm W(\mathcal Z)
=
2L\sum_{s\in\mathcal Z_{+}}
\Bigl(\Re(s)-\frac12\Bigr).
\]
由于每一项都不小于 \(\delta_{\min}(\mathcal Z)\)，便有
\[
\mathrm W(\mathcal Z)\ge 2L\,\mathrm{Ind}(\mathcal Z)\,\delta_{\min}(\mathcal Z),
\]
即
\[
\mathrm{Ind}(\mathcal Z)
\le
\frac{\mathrm W(\mathcal Z)}{2L\,\delta_{\min}(\mathcal Z)}.
\]
再由定理 \ref{thm:xi-offslice-blaschke-origin-modulus}，
\[
\mathrm W(\mathcal Z)=-\log|B_{\mathcal Z}(0)|,
\]
从而得到第一条闭式。另一方面，命题 \ref{prop:xi-offslice-model-space-dimension} 给出
\[
\dim K_{B_{\mathcal Z}}=\mathrm{Ind}(\mathcal Z),
\]
遂得第二条闭式。

等号情形同样直接：上式中的和达到下界，当且仅当每个求和项都等于 \(\delta_{\min}(\mathcal Z)\)，亦即全部 \(\Re(s)-\tfrac12\) 完全相同。证毕。
\end{proof}

\begin{theorem}[Blaschke 原点模在 strip 约束下的残差维数下界]\label{thm:xi-time-part65e-strip-blaschke-origin-modulus-index-lower-bound}
\leanverified{paper\_xi\_time\_part65e\_strip\_blaschke\_origin\_modulus\_index\_lower\_bound}
沿用定义 \ref{def:xi-offslice-ind-weight} 的记号，并设 \(\mathcal Z_{+}\neq\varnothing\) 为有限离切片零点多重集，\(B_{\mathcal Z}\) 为其对应的有限 Blaschke 乘积。
若存在常数 \(\eta>0\) 使
\[
0<\Re(s)-\frac12\le \eta
\qquad(s\in\mathcal Z_{+}),
\]
则有 sharp 下界
\[
\boxed{\;
\mathrm{Ind}(\mathcal Z)
\ge
\frac{\mathrm W(\mathcal Z)}{2L\,\eta}
=
\frac{-\log|B_{\mathcal Z}(0)|}{2L\,\eta}.\;}
\]
等价地，
\[
\boxed{\;
\dim K_{B_{\mathcal Z}}
\ge
\frac{-\log|B_{\mathcal Z}(0)|}{2L\,\eta}.\;}
\]
并且等号成立当且仅当全部离切片零点具有相同的实部偏差
\[
\Re(s)-\frac12\equiv \eta
\qquad(s\in\mathcal Z_{+}).
\]
\end{theorem}
\begin{proof}
由定理 \ref{thm:xi-offslice-realpart-sum-law}，
\[
\mathrm W(\mathcal Z)
=
2L\sum_{s\in\mathcal Z_{+}}
\Bigl(\Re(s)-\frac12\Bigr).
\]
由于每一项都不超过 \(\eta\)，故
\[
\mathrm W(\mathcal Z)\le 2L\,\eta\,\mathrm{Ind}(\mathcal Z),
\]
从而
\[
\mathrm{Ind}(\mathcal Z)\ge \frac{\mathrm W(\mathcal Z)}{2L\,\eta}.
\]
再由定理 \ref{thm:xi-offslice-blaschke-origin-modulus}，
\(
\mathrm W(\mathcal Z)=-\log|B_{\mathcal Z}(0)|
\)，
即得第一条闭式。
另一方面，由命题 \ref{prop:xi-offslice-model-space-dimension}，
\[
\dim K_{B_{\mathcal Z}}=\mathrm{Ind}(\mathcal Z),
\]
于是第二条闭式随即成立。

等号成立当且仅当上式求和中的每一项都达到其上界 \(\eta\)，亦即全部 \(\Re(s)-\tfrac12\) 完全相同且等于 \(\eta\)。证毕。
\end{proof}

\begin{theorem}[Rouch\'e 边界裕度所给出的 RH 有限证书]\label{thm:xi-time-part65e-rouche-boundary-margin-rh-certificate}
沿用 \eqref{eq:xi-ramanujan-horizon-Fzeta} 的圆盘完成化函数 \(F_{\zeta}\)。设 \(\{F_N\}_{N\ge 1}\) 为一列解析函数 \(F_N:\mathbb{D}\to\CC\)，并满足定理 \ref{thm:xi-hurwitz-horizon-reduction} 的自对偶完成化一致性。若存在一列半径
\[
\rho_N\uparrow 1
\qquad(N\to\infty)
\]
使得对每个 \(N\) 都有
\[
F_N\text{ 在 }\overline{D(0,\rho_N)}\text{ 内无零点}
\]
且
\[
\sup_{|w|=\rho_N}|F_{\zeta}(w)-F_N(w)|
<
\inf_{|w|=\rho_N}|F_N(w)|,
\]
则 RH 成立。
\end{theorem}
\begin{proof}
固定 \(N\)。由于 \(F_N\) 在 \(\overline{D(0,\rho_N)}\) 内无零点，故
\[
\inf_{|w|=\rho_N}|F_N(w)|>0.
\]
于是由 Rouch\'e 定理，\(F_{\zeta}\) 与 \(F_N\) 在圆盘 \(D(0,\rho_N)\) 内具有相同的零点总数（按重数计）。而 \(F_N\) 在该圆盘内无零点，故 \(F_{\zeta}\) 在 \(D(0,\rho_N)\) 内亦无零点。

由于 \(\rho_N\uparrow 1\)，对任意 \(0<\rho<1\) 都可取 \(N\) 充分大使 \(\rho<\rho_N\)。于是 \(F_{\zeta}\) 在 \(D(0,\rho)\) 内无零点。因为这对每个 \(\rho<1\) 都成立，便知 \(F_{\zeta}\) 在整个 \(\mathbb D\) 内无零点。最后由定理 \ref{thm:app-rh-iff-disk-zero-free}，RH 成立。证毕。
\end{proof}

\begin{theorem}[window-\(6\) 的几何可实现 boundary parity 子群恰为对角 \texorpdfstring{$\ZZ_2$}{Z2}，且精确损失两比特]\label{thm:xi-time-part65e-window6-geometric-diagonal-z2-parity}
\leanverified{paper\_xi\_time\_part65e\_window6\_geometric\_diagonal\_z2\_parity}
记
\[
P_6:=C_{\mathrm{bdry}}\cong (\ZZ/2\ZZ)^3
\]
为命题 \ref{prop:xi-time-part60ab3-window6-boundary-center-residual} 的 boundary parity 中心子群。若只允许保持稳定类型标签且由 \(\Aut(Q_6)\) 实现的强局域 geometric uplift，则其几何可实现的 boundary parity 翻转子群恰为
\[
P_6^{\mathrm{geo}}=\langle\zeta\rangle\cong \ZZ/2\ZZ,
\]
其中 \(\zeta\) 为定理 \ref{thm:xi-time-part53da-boundary-z6-torsor} 中同时翻转三条两点纤维的全局 sheet-flip。更精确地，在坐标识别
\[
P_6\cong (\ZZ/2\ZZ)^3
\]
下有
\[
P_6^{\mathrm{geo}}=\langle(1,1,1)\rangle.
\]
因此
\[
P_6/P_6^{\mathrm{geo}}\cong (\ZZ/2\ZZ)^2.
\]
换言之，局域几何可实现性仅保留三位 boundary parity 的同步对角方向，并精确剥离两个独立的二值自由度。
\end{theorem}
\begin{proof}
由命题 \ref{prop:xi-time-part60ab3-window6-boundary-center-residual}，三条 boundary 二点纤维各自产生一个彼此独立的中心 sheet-flip，故
\[
P_6\cong (\ZZ/2\ZZ)^3.
\]
另一方面，由定理 \ref{thm:xi-time-part53da-boundary-z6-torsor}，\(\widetilde X_6^{\mathrm{bdry}}\) 上存在一个全局对合
\[
\zeta:\widetilde X_6^{\mathrm{bdry}}\to \widetilde X_6^{\mathrm{bdry}},
\]
它在每条两点纤维上都施加唯一非平凡交换，因此在上述三维 parity 坐标中正对应于同步翻转 \((1,1,1)\)。

再由结论 \ref{con:xi-terminal-window6-local-uplift-so10-unique}，在保持稳定类型标签 \(F_6\) 不变且由 \(\Aut(Q_6)\) 实现的强局域 uplift 口径下，除平凡位移外仅允许 \(\delta=\pm 34\)。而在 window-\(6\) 的 boundary uplift 中，同一纤维内两点之差正恒等于 \(34\)（定理 \ref{thm:xi-time-part53da-boundary-z6-torsor} 的设定），故这唯一允许的几何非平凡分支恰好就是 \(\zeta\)。因此
\[
P_6^{\mathrm{geo}}=\langle \zeta\rangle=\langle(1,1,1)\rangle\cong \ZZ/2\ZZ.
\]
最后，\((\ZZ/2\ZZ)^3\) 除以其一维对角子群后得到
\[
(\ZZ/2\ZZ)^3/\langle(1,1,1)\rangle\cong (\ZZ/2\ZZ)^2,
\]
从而 cokernel 精确具有两位二进自由度。证毕。
\end{proof}

\begin{theorem}[window-\(6\) boundary parity 的无协议门跨尺度 \(2\)-挠中心延拓不存在]\label{thm:xi-time-part65e-window6-no-gateless-z2-central-extension}
\leanverified{paper\_xi\_time\_part65e\_window6\_no\_gateless\_z2\_central\_extension}
不存在某个固定紧连通群 \(G\) 与一族仅依赖于分辨率的规则
\[
\mathfrak F_m:X_m^{\mathrm{bdry}}\longrightarrow Z(G)[2]
\qquad(m\in\{6,7,8\}),
\]
其中
\[
Z(G)[2]:=\{z\in Z(G):z^2=e\},
\]
使得同时满足下列三条条件：
\begin{enumerate}
  \item 在 \(m=6\) 时，\(\mathfrak F_6\) 恢复推论 \ref{cor:terminal-delta-parity-rule} 中三条 boundary 纤维所给出的独立 \((\ZZ_2)^3\)-sheet parity；
  \item 该规则对 dyadic uplift 的 boundary 纤维置换保持自然；
  \item 构造过程中不引入任何额外协议门或额外读出寄存器。
\end{enumerate}
\end{theorem}
\begin{proof}
反设存在这样一族规则。条件 (1) 表明：在 \(m=6\) 的 boundary 扇区上，每条二点纤维都被赋予一个非平凡的二值中心选择；这正是推论 \ref{cor:terminal-delta-parity-rule} 中由
\[
\Delta_6^{\mathrm{bdry}}=\{0,34\}
\]
所诱导的 sheet parity。若条件 (2) 与 (3) 也同时成立，则这一二值对象必须沿 dyadic uplift 在更高分辨率的 boundary 纤维上继续以“置换自然且无需外加侧信息”的方式存在。

然而，命题 \ref{prop:terminal-sign-covariant-label-boundary-bifurcation} 已经给出严格判据：这样的符号协变二值对象存在当且仅当纤维基数等于 \(2\)。同一命题并明确指出，在 boundary uplift 口径下有
\[
m=6:\ \#(\mathrm{Fold}^{\mathrm{bin}}_6)^{-1}(w)=2,
\]
\[
m=7,8:\ \#(\mathrm{Fold}^{\mathrm{bin}}_m)^{-1}(w)=3,
\qquad
\Delta_7^{\mathrm{bdry}}=\{0,55,89\},
\qquad
\Delta_8^{\mathrm{bdry}}=\{0,89,144\}.
\]
因此 \(m=7,8\) 的三点 boundary 纤维上根本不存在所需的置换自然二值延拓，这与条件 (2)--(3) 矛盾。故这样的统一规则不存在。

最后，备注 \ref{rem:window6-sheet-parity-nonpersistent} 已指出：若欲把 window-\(6\) 的 \(\ZZ_2\) sheet parity 提升为跨尺度残余选择律，就必须额外施加可审计协议门以重新锁定二层覆盖，或引入额外读出寄存器。并且定理 \ref{thm:xi-boundary-uplift-mixed-radix-register-lower-bound} 进一步量化了这一下限：若要求完整继承 window-\(6\) 的三位 sheet parity，则必有
\[
r_{\mathrm{ext}}(7)\ge 11,
\qquad
r_{\mathrm{ext}}(8)\ge 16.
\]
证毕。
\end{proof}

\noindent
由此，Toeplitz 主子式、离切片 Blaschke 数据、圆盘完成化边界比较与 window-\(6\) boundary uplift 这五条原本分属不同模块的链条，都被进一步压缩为有限局域证书：Toeplitz 可实现性不再需要搜索整个正定锥，而可在对数坐标下完全转写为离散凹锥；Blaschke 原点模不再只是总偏离负载，而在给定最小偏离阈值时直接控制残差维数；RH 的一个充分条件不再需要整列紧集一致逼近，而只需一列逼近边界的圆周裕度不等式；至于 window-\(6\) 的 \((\ZZ_2)^3\) boundary parity，则在局域几何 uplift 口径下仅保留同步对角 \(\ZZ_2\) 方向，而其跨尺度延拓仍必须显式支付新的协议成本。

\endinput
